% Abstract draft from docs/11 §1. Bracketed claims stay hypotheses until bench/results/ exists.
Large language model (LLM) agents now name the faulty service in most cases of some public root-cause analysis (RCA) benchmarks, yet stricter tasks, unseen failures and planted decoys show that their answers cannot be used without checking.
We present \sys, an open RCA engine for microservices in which statistical rankers and a local open-weight LLM only \emph{propose} hypotheses and a deterministic verifier \emph{decides}.
Every claim the LLM makes is a named predicate that code re-evaluates on stored telemetry for every candidate, so the LLM can surface discriminating evidence but cannot argue for its favourite.
The verifier returns a split-conformal candidate set with a target coverage of 90\%, may answer \emph{no fault} or \emph{inconclusive}, and seals each verdict in a receipt from which anyone can replay it to the same hash.
We hypothesise that (i) the deterministic pipeline is not worse than the best simple baseline on any RCAEval RE2 system, (ii) code-verified proposals from a local model of at most 10B parameters do not lower accuracy and resist planted decoys, (iii) the candidate sets reach nominal coverage within a benchmark family, and (iv) the design is robust to alert-time shifts, missing telemetry and decoy change events.
We evaluate per system with oracle and detected alert times, against chance, N-Sigma, BARO and an alert-count baseline, using clustered paired statistics, and we release all per-case outputs.
\leadq{headline numbers for RQ1--RQ6 go here from bench/results/ at M4.}
\sys runs on a 16\,GB laptop in \resnum{re2-oracle-rq1-v1}{p95 wall time per case, veyltrace:off} seconds per case without a GPU.
